\documentclass[11pt]{article}
\usepackage[margin=1in]{geometry}
\usepackage{amsmath,amssymb,amsthm,booktabs,xcolor,hyperref}
\hypersetup{colorlinks=true,allcolors=blue!45!black}
\newtheorem{theorem}{Theorem}
\title{Thirty-eight Prime Outputs Refute the Ova Counting Bound\\A Complete Formal Counterexample to MathDB 350342}
\author{Myeongjin Kang\\\small AI-assisted verification note}
\date{October 25, 2026}
\begin{document}
\maketitle
\begin{abstract}
We verify a finite counterexample to the proposed upper bound on the number of primes produced by the Ova-Ova-Prime-Ova construction. The input primes 677 and 401 produce at least 38 distinct prime outputs. This disproves both the bound 10 in MathDB 350342 and the bound 13 in the original 2020 publication. Every displayed primality claim, residue condition, and cardinality inequality is checked in Lean. The witness was posted previously by Shivam Patel; this note claims verification, not discovery.
\end{abstract}

\section{The construction and the source discrepancy}
For a natural number $p$, write
\[
 k_p=p\bmod360,\qquad f_p=\lfloor p/360\rfloor,
 \qquad p=k_p+360f_p.
\]
For two prime inputs $p,q$, the construction in Theorem 7 of \cite{A2020} considers prime residue classes $r$ modulo 360 for which
\[
 c_r=k_p+k_q+2-r
\]
is prime. Whenever
\[
 v_r=r+360(f_p+f_q)
\]
is also prime, $v_r$ is a resulting prime. Here $1\leq r<360$ and $r\leq k_p+k_q+2$, so the subtraction is an ordinary nonnegative integer subtraction. The natural meaning of the number of resulting primes is the number of distinct such values $v_r$.

Conjecture 1 of the original paper \cite{A2020}, printed on page 117, proposes an upper bound of 13. The later arXiv paper \cite{A2021} cites Theorem 7 and reuses the construction in its Goldbach discussion, but the retrieved versions do not supply the claimed bound 10. MathDB 350342 explicitly records the bound 10 \cite{DB}. These are different source statements. We do not attribute the number 10 to the original Conjecture 1. The counterexample below refutes both bounds, so this bibliographic discrepancy does not affect the mathematical outcome.

\begin{theorem}\label{thm:main}
For the input primes $p=677$ and $q=401$, the construction has at least 38 distinct prime outputs. In particular its output count exceeds both 10 and 13.
\end{theorem}

\section{The witness and its complete verification}
The inputs satisfy
\[
 677=317+360\cdot1,\qquad401=41+360\cdot1.
\]
Both are prime. For these inputs the two output conditions simplify to
\[
 c_r=360-r\text{ is prime},\qquad v_r=720+r\text{ is prime}.
\]
The following table provides 38 residues and the corresponding prime pairs.
\begin{center}\small
\begin{tabular}{rrr@{\hspace{1cm}}rrr}
\toprule
$r$ & $360-r$ & $720+r$ & $r$ & $360-r$ & $720+r$\\
\midrule
7&353&727&221&139&941\\
13&347&733&233&127&953\\
23&337&743&247&113&967\\
49&311&769&251&109&971\\
53&307&773&257&103&977\\
67&293&787&263&97&983\\
77&283&797&271&89&991\\
89&271&809&277&83&997\\
91&269&811&289&71&1009\\
103&257&823&293&67&1013\\
109&251&829&299&61&1019\\
119&241&839&301&59&1021\\
133&227&853&313&47&1033\\
137&223&857&319&41&1039\\
161&199&881&329&31&1049\\
163&197&883&331&29&1051\\
167&193&887&341&19&1061\\
187&173&907&343&17&1063\\
209&151&929&349&11&1069\\
\bottomrule
\end{tabular}
\end{center}
All residues lie strictly between 0 and 360. Every integer in either prime column is prime. These finite assertions can be checked by trial division using the primes at most 31, since the largest entry is 1069 and $33^2>1069$. Lean's numerical normalization produces checked proofs of the primality assertions rather than accepting a primality oracle.

Each residue also genuinely represents a prime residue class: its output $v_r$ is prime and
\[
 v_r\bmod360=(720+r)\bmod360=r.
\]
Thus the residue-class requirement in the original construction has not been replaced with a heuristic coprimality test or omitted. The prime $v_r$ itself is an explicit witness to that requirement.

Finally $r\mapsto720+r$ is injective. Since the table lists 38 different residues, it lists 38 different output primes. This proves Theorem~\ref{thm:main}. Exhaustively ruling out further outputs is unnecessary: a lower bound of 38 already contradicts either universal upper bound.

\section{A precise finite-set formulation}
To make the full bridge to the source explicit, let $E(p,q,r)$ mean the conjunction
\begin{align*}
 &0<r<360,\
 &r\leq k_p+k_q+2,\
 &\operatorname{Prime}(k_p+k_q+2-r),\
 &\operatorname{Prime}(r+360(f_p+f_q)).
\end{align*}
Define finite sets
\begin{align*}
 R(p,q)&=\{r\in\{0,\ldots,359\}:E(p,q,r)\},\\
 O(p,q)&=\{r+360(f_p+f_q):r\in R(p,q)\}.
\end{align*}
Every member of $O(p,q)$ is prime. For every $r\in R(p,q)$, its output prime is a witness that $r$ belongs to the original set of prime residue classes. Conversely, a residue satisfying the original source's output conditions satisfies $E$. Hence the finite set $O(p,q)$ implements the original construction's resulting primes, rather than a larger relaxed set.

The identity
\[
 (k_p+k_q+2-r)+(r+360(f_p+f_q))=p+q+2
\]
holds whenever $E(p,q,r)$ holds. This confirms the source decomposition as well as the output conditions. The formal proof includes this identity for arbitrary natural inputs and every admissible residue.

The finite set $W$ of table residues satisfies $|W|=38$ and $W\subseteq R(677,401)$. Therefore
\[
 38=|W|\leq|R(677,401)|=|O(677,401)|.
\]
The last equality follows from injectivity of addition of the fixed frequency term. This is the exact cardinality argument checked by Lean.

\section{Formal coverage, attribution, and scope}
The module \texttt{formal/OvaCount.lean} defines the input residue and frequency functions, the full numerical eligibility predicate, and the two finite sets above. Its declarations verify:
\begin{itemize}
\item primality of every output and the source decomposition for all inputs and admissible residues;
\item equality of residue and output cardinalities, using injectivity;
\item primality and residue/frequency values of the two actual inputs;
\item cardinality 38 of the displayed witness set, and eligibility of every member;
\item an explicit prime witness to the prime-residue-class requirement;
\item existence of two prime inputs with at least 38 outputs, hence more than 10 and more than 13.
\end{itemize}
The definition \texttt{SourceEligible} explicitly conjoins the prime-residue-class condition with numerical eligibility. The equivalence \texttt{source\_condition\_iff} and the finite-set identity \texttt{source\_outputs\_eq} prove that this explicit source formulation counts the same outputs.

The final theorem \texttt{OvaCount.prime\_count\_counterexample} uses this source-formulated output set. A finite universal statement over the witness set is proved by splitting the 38 possible residues and checking the arithmetic and primality goals. The small finite-set cardinality is proved by kernel reduction. No unchecked native computation, admitted claim, or extra mathematical axiom is used.

The pinned Lean 4.34.1 compiler checked every local declaration with warnings as errors. The complete axiom audit uses only propositional extensionality, classical choice, and quotient soundness. All dependencies, proof hashes, and source hashes are recorded alongside the module. The manuscript and source interpretation await the parent campaign's independent review.

The same inputs and the same 38 residues were already supplied by Shivam Patel in the MathDB response dated August 20, 2026. That prior response claims an exact total of 38. Our formal theorem deliberately asserts only the lower bound needed for the full disproof. The lower bound is exact enough to resolve the conjecture; it is not a partial result on a still-open case. We claim neither the witness's discovery nor a proof or disproof of the source's separate general at-least-one-output theorem. No Goldbach assertion is used anywhere in this argument.

\begin{thebibliography}{9}
\bibitem{A2020} Y. A. Acevedo Agudelo, \emph{A complete classification of the Mersenne's primes and its implications for computing}, Revista Polit\'ecnica 16(32) (2020), 111--119, Theorem 7 and Conjecture 1. \url{https://doi.org/10.33571/rpolitec.v16n32a10}.
\bibitem{A2021} Y. A. Acevedo Agudelo, \emph{Prime numbers. An alternative study using ova-angular rotations}, arXiv:2104.04522v1 and v2 (2021), discussion citing Theorem 7 of the preceding reference. \url{https://arxiv.org/abs/2104.04522}.
\bibitem{DB} MathDB, post 350342, \emph{The upper-bound conjecture for the Ova-Ova-Prime-Ova construction}; counterexample by Shivam Patel, August 20, 2026. \url{https://mathdb.com/p/350342}.
\end{thebibliography}
\end{document}
